% =============================================================================
% APPENDIX CCC: METHOD-DOCTYPE: 8D Document Type Clustering
% =============================================================================
% Module: BCM2_METHOD_DOCTYPE
% Category: METHOD
% Version: 1.0 (January 2026)
% Status: Draft
% Language: English
%
% Computational framework for characterizing and clustering document types
% in an 8-dimensional audience-centered space. Enables automated document
% type identification, gap analysis, and style recommendations.
%
% =============================================================================

% -----------------------------------------------------------------------------
% HEADER BLOCK
% -----------------------------------------------------------------------------
% APPENDIX CCC: METHOD-DOCTYPE: 8D Document Type Clustering
% Module: BCM2_METHOD_DOCTYPE
% Version: 1.0 (January 2026)
% Status: Draft
%
% This appendix provides a computational framework for characterizing document
% types as points in an 8-dimensional audience-centered space, enabling
% systematic clustering, gap analysis, and style recommendations.
% -----------------------------------------------------------------------------

% =============================================================================
% CROSS-REFERENCE STRUCTURE
% =============================================================================
%
% CHAPTER LINKAGE (Required):
% - Primary Chapter: Chapter 17: Policy Implications (document design)
% - Secondary Chapters: Chapter 4x: Calibration (LLM measurement)
%
% DEPENDENCIES (what this appendix REQUIRES):
% - Appendix LLM1 (METHOD-LLMMC): LLM-based measurement methodology
% - Appendix UNMAPPED_GLS (REF-GLOSSARY): Standard terminology
%
% DEPENDENTS (what REQUIRES this appendix):
% - Quality management system: Document style validation
% - Appendix UNMAPPED_DTE (REF-DOCTYPE): Computational complement to theory
%
% =============================================================================

\begin{tcolorbox}[colback=purple!5!white, colframe=purple!75!black,
    title=Chapter Linkage]

\textbf{Primary Chapter:} Chapter 17: Policy Implications
\begin{itemize}[nosep]
\item Section 17.3: Communication Design $\leftrightarrow$ This Appendix Section 3
\item Section 17.4: Stakeholder Analysis $\leftrightarrow$ This Appendix Section 2
\end{itemize}

\textbf{Secondary Chapters:}
\begin{itemize}[nosep]
\item Chapter 4x: Calibration --- LLM-based measurement of D1-D8
\item Chapter 9: Context --- $\Psi$ parallels to document context
\end{itemize}

\textbf{Reading Guide:}
\begin{itemize}[nosep]
\item For conceptual overview: Read Section 1 (Fundamental Question)
\item For implementation: See Section 3 (Computational Framework)
\item For applications: See Section 4 (Worked Examples)
\end{itemize}

\end{tcolorbox}

\chapter{METHOD-DOCTYPE: 8D Document Type Clustering}
\label{app:doctype}

\begin{abstract}
This appendix presents a computational framework for characterizing document types
as points in an 8-dimensional space. The eight dimensions---Knowledge, Proximity,
Scope, Time, Goal, Context, Emotion, and Persistence---enable automated clustering
of document types, identification of gaps in the document space, and systematic
style recommendations. The framework treats document types not as fixed categories
but as emergent clusters in a continuous audience-centered space.
\end{abstract}

\begin{tcolorbox}[colback=blue!5!white, colframe=blue!75!black,
    title=Quick Reference: 8D Document Type Framework]
\textbf{Appendix:} DTC (METHOD)


\textbf{Core Question:} How can we systematically characterize and cluster
all document types in a unified framework?

\textbf{Key Formula:}
\begin{equation}
\text{DocType}(Z) = f(D_1, D_2, D_3, D_4, D_5, D_6, D_7, D_8) \in [0,1]^8
\end{equation}

\textbf{Key Concepts:}
\begin{itemize}[nosep]
\item \textbf{8D Space}: Documents characterized by 8 audience-centered dimensions
\item \textbf{Emergent Clusters}: Document categories emerge from dimension clustering
\item \textbf{LLM Measurement}: Dimensions measured via LLM analysis
\end{itemize}

\textbf{Cross-References:} METHOD-LLMMC (AN), REF-GLOSSARY (G), CORE-WHEN (V)
\end{tcolorbox}

% -----------------------------------------------------------------------------
% APPENDIX SCOPE BOX
% -----------------------------------------------------------------------------

\begin{tcolorbox}[
    colback=orange!5!white,
    colframe=orange!75!black,
    title={\textbf{Appendix Scope: Ziel / In-Scope / Out-of-Scope / Constraints}},
    fonttitle=\bfseries
]

\textbf{Ziel (Objective):}
\begin{itemize}[nosep]
    \item Core Question:
\end{itemize}

\textbf{In-Scope:}
\begin{itemize}[nosep]
    \item The Fundamental Question
    \item Core Theory: The 8D Document Space
    \item Axioms and Formal Definitions
    \item Key Results
    \item Computational Framework
\end{itemize}

\textbf{Out-of-Scope (delegiert an andere Appendices):}
\begin{itemize}[nosep]
    \item REF-GLOSSARY $\rightarrow$ Appendix UNMAPPED_GLS
    \item REF-DOCTYPE $\rightarrow$ Appendix UNMAPPED_DTE
    \item Central Glossary $\rightarrow$ Appendix UNMAPPED_GLS
    \item CORE-WHEN $\rightarrow$ Appendix CTW
    \item REF-DOCTYPE --- theoretical foundations $\rightarrow$ Appendix UNMAPPED_DTE
\end{itemize}

\textbf{Constraints (Anwendungsgrenzen):}
\begin{itemize}[nosep]
    \item [TODO: Define when this appendix does NOT apply]
    \item [TODO: Define required prerequisites]
    \item [TODO: Define known limitations]
\end{itemize}

\textbf{Lieferobjekte (Deliverables):}
\begin{itemize}[nosep]
    \item 11 Table(s)
    \item 7 Equation(s)
    \item 9 Axiom(s)
    \item Worked Example(s)
\end{itemize}

\end{tcolorbox}




%%%%%%%%%%%%%%%%%%%%%%%%%%%%%%%%%%%%%%%%%%%%%%%%%%%%%%%%%%%%%%%%%%%%%%%%%%%%%%%
\section{The Fundamental Question}
\label{sec:doctype-fundamental}
%%%%%%%%%%%%%%%%%%%%%%%%%%%%%%%%%%%%%%%%%%%%%%%%%%%%%%%%%%%%%%%%%%%%%%%%%%%%%%%

\subsection{Why This Matters}

Document types are typically treated as fixed, discrete categories: ``journal article,''
``policy brief,'' ``novel.'' This categorical approach has limitations:

\begin{enumerate}
\item \textbf{Incompleteness}: New document types emerge constantly (tweets, podcasts, etc.)
\item \textbf{Ambiguity}: Hybrid documents don't fit discrete categories
\item \textbf{No Guidance}: Categories don't tell authors \textit{how} to write
\item \textbf{Culture-Bound}: Categories vary across disciplines and cultures
\end{enumerate}

\subsection{The Core Problem}

We need a framework that can:
\begin{itemize}[nosep]
\item Characterize \textit{any} document type (existing or novel)
\item Identify natural clusters without predefined categories
\item Guide document design based on audience characteristics
\item Enable computational analysis and recommendation
\end{itemize}

\begin{tcolorbox}[colback=red!5!white, colframe=red!75!black,
    title=The Fundamental Question]
\textbf{How can we represent all document types as points in a continuous,
audience-centered space that enables systematic characterization, clustering,
and style recommendation?}
\end{tcolorbox}


%%%%%%%%%%%%%%%%%%%%%%%%%%%%%%%%%%%%%%%%%%%%%%%%%%%%%%%%%%%%%%%%%%%%%%%%%%%%%%%
\section{Core Theory: The 8D Document Space}
\label{sec:doctype-theory}
%%%%%%%%%%%%%%%%%%%%%%%%%%%%%%%%%%%%%%%%%%%%%%%%%%%%%%%%%%%%%%%%%%%%%%%%%%%%%%%

\subsection{The Eight Dimensions}

Every document can be characterized by its position on eight dimensions,
each reflecting a property of the target audience or communication context:

\begin{table}[htbp]
\centering
\caption{The 8D Document Type Dimensions}
\label{tab:8d-dimensions}
\renewcommand{\arraystretch}{1.4}
\begin{tabular}{@{}clll@{}}
\toprule
\textbf{D} & \textbf{Dimension} & \textbf{Low (0)} & \textbf{High (1)} \\
\midrule
$D_1$ & \textbf{Wissen} (Knowledge) & Layperson & Expert \\
$D_2$ & \textbf{N\"ahe} (Proximity) & Distant field & Same discipline \\
$D_3$ & \textbf{Reichweite} (Scope) & Personal & Societal \\
$D_4$ & \textbf{Zeit} (Time) & Little reading time & Much reading time \\
$D_5$ & \textbf{Ziel} (Goal) & See subcategories & See subcategories \\
$D_6$ & \textbf{Kontext} (Context) & Internal & External \\
$D_7$ & \textbf{Emotion} & Sachlich (factual) & Emotional \\
$D_8$ & \textbf{Persistenz} (Persistence) & Ephemeral & Archival \\
\bottomrule
\end{tabular}
\end{table}

\subsection{D5: The Goal Dimension (Expanded)}

Dimension 5 (Goal) requires special treatment as it is categorical rather than
continuous. We identify seven primary communication goals:

\begin{table}[htbp]
\centering
\caption{D5 Goal Categories}
\label{tab:d5-goals}
\renewcommand{\arraystretch}{1.3}
\begin{tabular}{@{}cll@{}}
\toprule
\textbf{Code} & \textbf{Goal} & \textbf{Primary Intent} \\
\midrule
$G_1$ & Informieren (Inform) & Transfer knowledge \\
$G_2$ & Handeln (Act) & Trigger specific action \\
$G_3$ & \"Uberzeugen (Persuade) & Change beliefs/attitudes \\
$G_4$ & Unterhalten (Entertain) & Provide enjoyment \\
$G_5$ & Erleben (Experience) & Create aesthetic experience \\
$G_6$ & Verbinden (Connect) & Build relationship \\
$G_7$ & Ausdr\"ucken (Express) & Convey emotion/opinion \\
\bottomrule
\end{tabular}
\end{table}

For computational purposes, $D_5$ can be represented as a 7-dimensional one-hot
vector or as the dominant goal index.

\subsection{The Document Type Function}

\begin{tcolorbox}[colback=gray!5!white, colframe=gray!75!black,
    title=Definition: Document Type Vector]

A document type $\mathcal{D}$ is characterized by its 8D vector:

\begin{equation}
\mathcal{D} = (D_1, D_2, D_3, D_4, D_5, D_6, D_7, D_8)
\end{equation}

where:
\begin{itemize}[nosep]
\item $D_1, D_2, D_3, D_4, D_6, D_7, D_8 \in [0, 1]$ (continuous)
\item $D_5 \in \{G_1, G_2, G_3, G_4, G_5, G_6, G_7\}$ (categorical)
\end{itemize}

\end{tcolorbox}


%%%%%%%%%%%%%%%%%%%%%%%%%%%%%%%%%%%%%%%%%%%%%%%%%%%%%%%%%%%%%%%%%%%%%%%%%%%%%%%
\section{Axioms and Formal Definitions}
\label{sec:doctype-axioms}
%%%%%%%%%%%%%%%%%%%%%%%%%%%%%%%%%%%%%%%%%%%%%%%%%%%%%%%%%%%%%%%%%%%%%%%%%%%%%%%

\begin{tcolorbox}[colback=gray!5!white, colframe=gray!75!black,
    title=Axiom DOC-1: Dimension Completeness]
\textbf{Statement:} Every document can be characterized by its position in the 8D space.

\begin{equation}
\forall d \in \mathcal{D}_{all}: \exists (D_1, ..., D_8) \text{ such that } d = f(D_1, ..., D_8)
\end{equation}

\textbf{Interpretation:} The 8 dimensions are sufficient to characterize any document type.

\textbf{Implication:} No additional dimensions are required for document classification.
\end{tcolorbox}

\begin{tcolorbox}[colback=gray!5!white, colframe=gray!75!black,
    title=Axiom DOC-2: Dimension Independence]
\textbf{Statement:} The 8 dimensions are mutually independent.

\begin{equation}
\text{Cov}(D_i, D_j) = 0 \quad \forall i \neq j
\end{equation}

\textbf{Interpretation:} Knowledge of one dimension provides no information about others.

\textbf{Implication:} Each dimension captures unique variance in document types.
\end{tcolorbox}

\begin{tcolorbox}[colback=gray!5!white, colframe=gray!75!black,
    title=Axiom DOC-3: Cluster Emergence]
\textbf{Statement:} Document type categories emerge from density in the 8D space.

\begin{equation}
\text{Category}(d) = \arg\max_C P(d \in C | D_1, ..., D_8)
\end{equation}

\textbf{Interpretation:} Categories are not predefined but emerge from clustering.

\textbf{Implication:} New document types naturally find their place in the space.
\end{tcolorbox}


%%%%%%%%%%%%%%%%%%%%%%%%%%%%%%%%%%%%%%%%%%%%%%%%%%%%%%%%%%%%%%%%%%%%%%%%%%%%%%%
\section{Key Results}
\label{sec:doctype-results}
%%%%%%%%%%%%%%%%%%%%%%%%%%%%%%%%%%%%%%%%%%%%%%%%%%%%%%%%%%%%%%%%%%%%%%%%%%%%%%%

\begin{proposition}[Universality]
The 8D framework can represent any document type, including those not yet invented.
\end{proposition}

\begin{proof}
By Axiom DOC-1, every document has a position in $[0,1]^7 \times \{G_1,...,G_7\}$.
Since this space is continuous (except $D_5$), it can represent infinite variations.
New document types correspond to previously unoccupied regions of this space.
\end{proof}

\begin{proposition}[Cluster Stability]
Emergent clusters are stable under small perturbations of document coordinates.
\end{proposition}

\begin{proof}
Using DBSCAN with $\epsilon > 0$, cluster membership is determined by density.
Small perturbations $\delta D_i < \epsilon/\sqrt{8}$ cannot change cluster assignment
as the document remains within the $\epsilon$-neighborhood of its cluster core.
\end{proof}


%%%%%%%%%%%%%%%%%%%%%%%%%%%%%%%%%%%%%%%%%%%%%%%%%%%%%%%%%%%%%%%%%%%%%%%%%%%%%%%
\section{Computational Framework}
\label{sec:doctype-computation}
%%%%%%%%%%%%%%%%%%%%%%%%%%%%%%%%%%%%%%%%%%%%%%%%%%%%%%%%%%%%%%%%%%%%%%%%%%%%%%%

\subsection{Data Structure}

Documents are represented as vectors in Python:

\begin{verbatim}
import numpy as np
from dataclasses import dataclass
from enum import Enum
from typing import List, Dict, Tuple

class Goal(Enum):
    INFORM = 1      # Informieren
    ACT = 2         # Handeln
    PERSUADE = 3    # Ueberzeugen
    ENTERTAIN = 4   # Unterhalten
    EXPERIENCE = 5  # Erleben
    CONNECT = 6     # Verbinden
    EXPRESS = 7     # Ausdruecken

@dataclass
class DocumentType:
    """8D representation of a document type."""
    name: str
    d1_knowledge: float    # 0=layperson, 1=expert
    d2_proximity: float    # 0=distant, 1=same field
    d3_scope: float        # 0=personal, 1=societal
    d4_time: float         # 0=little, 1=much
    d5_goal: Goal          # categorical
    d6_context: float      # 0=internal, 1=external
    d7_emotion: float      # 0=factual, 1=emotional
    d8_persistence: float  # 0=ephemeral, 1=archival

    def to_vector(self) -> np.ndarray:
        """Convert to 8D numpy vector (goal as index)."""
        return np.array([
            self.d1_knowledge,
            self.d2_proximity,
            self.d3_scope,
            self.d4_time,
            self.d5_goal.value / 7.0,  # normalize to [0,1]
            self.d6_context,
            self.d7_emotion,
            self.d8_persistence
        ])
\end{verbatim}

\subsection{Example Document Types}

\begin{verbatim}
# Define example document types
DOCUMENT_TYPES = {
    "Journal Paper": DocumentType(
        name="Journal Paper",
        d1_knowledge=0.9,    # Expert audience
        d2_proximity=0.9,    # Same field
        d3_scope=0.7,        # Field-wide impact
        d4_time=0.8,         # Much reading time
        d5_goal=Goal.INFORM,
        d6_context=1.0,      # External (public)
        d7_emotion=0.1,      # Highly factual
        d8_persistence=0.95  # Archival
    ),
    "Tweet": DocumentType(
        name="Tweet",
        d1_knowledge=0.3,
        d2_proximity=0.2,
        d3_scope=0.5,
        d4_time=0.05,        # Seconds
        d5_goal=Goal.EXPRESS,
        d6_context=1.0,
        d7_emotion=0.6,
        d8_persistence=0.1   # Ephemeral
    ),
    "Love Letter": DocumentType(
        name="Love Letter",
        d1_knowledge=0.5,
        d2_proximity=0.1,
        d3_scope=0.1,        # Personal
        d4_time=0.7,
        d5_goal=Goal.CONNECT,
        d6_context=0.0,      # Internal
        d7_emotion=0.95,     # Highly emotional
        d8_persistence=0.6
    ),
    "Policy Brief": DocumentType(
        name="Policy Brief",
        d1_knowledge=0.6,
        d2_proximity=0.4,
        d3_scope=0.85,       # Societal impact
        d4_time=0.3,         # Quick read
        d5_goal=Goal.ACT,
        d6_context=1.0,
        d7_emotion=0.25,
        d8_persistence=0.7
    ),
    "Novel": DocumentType(
        name="Novel",
        d1_knowledge=0.3,
        d2_proximity=0.2,
        d3_scope=0.6,
        d4_time=0.95,        # Hours
        d5_goal=Goal.ENTERTAIN,
        d6_context=1.0,
        d7_emotion=0.8,
        d8_persistence=0.85
    ),
    "Internal Memo": DocumentType(
        name="Internal Memo",
        d1_knowledge=0.5,
        d2_proximity=0.8,
        d3_scope=0.3,        # Organizational
        d4_time=0.2,
        d5_goal=Goal.INFORM,
        d6_context=0.0,      # Internal
        d7_emotion=0.15,
        d8_persistence=0.4
    ),
}
\end{verbatim}

\subsection{Clustering Algorithm}

\begin{verbatim}
from sklearn.cluster import DBSCAN, KMeans
from sklearn.manifold import TSNE
from sklearn.preprocessing import StandardScaler
import matplotlib.pyplot as plt

def cluster_documents(
    documents: Dict[str, DocumentType],
    method: str = "dbscan",
    n_clusters: int = 5
) -> Tuple[np.ndarray, np.ndarray]:
    """
    Cluster document types in 8D space.

    Args:
        documents: Dictionary of DocumentType objects
        method: "dbscan" for density-based, "kmeans" for centroid-based
        n_clusters: Number of clusters (only for kmeans)

    Returns:
        labels: Cluster assignment for each document
        X: Feature matrix
    """
    # Convert to matrix
    names = list(documents.keys())
    X = np.array([documents[n].to_vector() for n in names])

    # Standardize features
    scaler = StandardScaler()
    X_scaled = scaler.fit_transform(X)

    # Cluster
    if method == "dbscan":
        clustering = DBSCAN(eps=0.5, min_samples=2).fit(X_scaled)
    else:
        clustering = KMeans(n_clusters=n_clusters).fit(X_scaled)

    return clustering.labels_, X_scaled, names

def visualize_clusters(
    X: np.ndarray,
    labels: np.ndarray,
    names: List[str]
) -> None:
    """Visualize 8D document space in 2D using t-SNE."""
    # Reduce to 2D
    tsne = TSNE(n_components=2, random_state=42, perplexity=min(5, len(X)-1))
    X_2d = tsne.fit_transform(X)

    # Plot
    plt.figure(figsize=(12, 8))
    scatter = plt.scatter(X_2d[:, 0], X_2d[:, 1],
                         c=labels, cmap='tab10', s=100)

    # Add labels
    for i, name in enumerate(names):
        plt.annotate(name, (X_2d[i, 0], X_2d[i, 1]),
                    fontsize=9, ha='center', va='bottom')

    plt.title("Document Types in 8D Space (t-SNE projection)")
    plt.colorbar(scatter, label="Cluster")
    plt.tight_layout()
    plt.savefig("document_clusters.png", dpi=150)
    plt.show()
\end{verbatim}

\subsection{Distance and Similarity}

\begin{verbatim}
def document_distance(
    doc1: DocumentType,
    doc2: DocumentType,
    weights: np.ndarray = None
) -> float:
    """
    Calculate weighted Euclidean distance between document types.

    Args:
        doc1, doc2: Document types to compare
        weights: Optional dimension weights (default: equal)

    Returns:
        Distance in [0, sqrt(8)] for unweighted
    """
    v1 = doc1.to_vector()
    v2 = doc2.to_vector()

    if weights is None:
        weights = np.ones(8)

    return np.sqrt(np.sum(weights * (v1 - v2)**2))

def find_nearest_type(
    target: DocumentType,
    corpus: Dict[str, DocumentType],
    n: int = 3
) -> List[Tuple[str, float]]:
    """Find n nearest document types to target."""
    distances = [
        (name, document_distance(target, doc))
        for name, doc in corpus.items()
    ]
    return sorted(distances, key=lambda x: x[1])[:n]
\end{verbatim}

\subsection{Gap Analysis}

\begin{verbatim}
def find_gaps(
    documents: Dict[str, DocumentType],
    resolution: int = 5
) -> List[np.ndarray]:
    """
    Find underrepresented regions in the 8D document space.

    Uses grid-based analysis to identify areas with no nearby
    document types.

    Args:
        documents: Known document types
        resolution: Grid points per dimension

    Returns:
        List of 8D coordinates representing gaps
    """
    X = np.array([d.to_vector() for d in documents.values()])

    gaps = []
    # Sample grid points (simplified - full grid is 5^8 = 390625 points)
    for _ in range(1000):
        # Random point in [0,1]^8
        point = np.random.rand(8)

        # Check distance to nearest document
        distances = np.linalg.norm(X - point, axis=1)
        if distances.min() > 0.5:  # Threshold for "gap"
            gaps.append(point)

    return gaps
\end{verbatim}


%%%%%%%%%%%%%%%%%%%%%%%%%%%%%%%%%%%%%%%%%%%%%%%%%%%%%%%%%%%%%%%%%%%%%%%%%%%%%%%
\section{LLM-Based Dimension Measurement}
\label{sec:doctype-llm}
%%%%%%%%%%%%%%%%%%%%%%%%%%%%%%%%%%%%%%%%%%%%%%%%%%%%%%%%%%%%%%%%%%%%%%%%%%%%%%%

Following the methodology in Appendix LLM1 (METHOD-LLMMC), we can use LLMs
to estimate the 8D coordinates of a document:

\begin{verbatim}
DIMENSION_PROMPT = """
Analyze this document and rate it on the following dimensions.
Return scores as JSON with values between 0.0 and 1.0.

Document: {document_text}

Dimensions to rate:
- d1_knowledge: How expert is the target audience?
  (0=layperson, 1=domain expert)
- d2_proximity: How close is the audience to the document's field?
  (0=completely unrelated, 1=same specialty)
- d3_scope: What is the impact scope?
  (0=personal/individual, 1=societal/global)
- d4_time: How much reading time is expected?
  (0=seconds, 1=hours)
- d5_goal: Primary communication goal (choose one):
  INFORM, ACT, PERSUADE, ENTERTAIN, EXPERIENCE, CONNECT, EXPRESS
- d6_context: Is this internal or external communication?
  (0=private/internal, 1=public/external)
- d7_emotion: How emotional vs factual is the tone?
  (0=purely factual, 1=highly emotional)
- d8_persistence: How long should this document last?
  (0=ephemeral, 1=archival/permanent)

Return JSON only:
"""

async def measure_document_8d(
    document_text: str,
    llm_client,
    n_samples: int = 5
) -> DocumentType:
    """
    Use LLM Monte Carlo to estimate 8D coordinates.

    Samples multiple LLM responses and aggregates for robustness.
    """
    samples = []

    for _ in range(n_samples):
        response = await llm_client.complete(
            DIMENSION_PROMPT.format(document_text=document_text[:2000])
        )
        try:
            scores = json.loads(response)
            samples.append(scores)
        except json.JSONDecodeError:
            continue

    # Aggregate (median for robustness)
    aggregated = {}
    for key in ['d1_knowledge', 'd2_proximity', 'd3_scope',
                'd4_time', 'd6_context', 'd7_emotion', 'd8_persistence']:
        values = [s[key] for s in samples if key in s]
        aggregated[key] = np.median(values) if values else 0.5

    # Mode for categorical goal
    goals = [s.get('d5_goal', 'INFORM') for s in samples]
    aggregated['d5_goal'] = max(set(goals), key=goals.count)

    return DocumentType(
        name="Measured Document",
        d1_knowledge=aggregated['d1_knowledge'],
        d2_proximity=aggregated['d2_proximity'],
        d3_scope=aggregated['d3_scope'],
        d4_time=aggregated['d4_time'],
        d5_goal=Goal[aggregated['d5_goal']],
        d6_context=aggregated['d6_context'],
        d7_emotion=aggregated['d7_emotion'],
        d8_persistence=aggregated['d8_persistence']
    )
\end{verbatim}


%%%%%%%%%%%%%%%%%%%%%%%%%%%%%%%%%%%%%%%%%%%%%%%%%%%%%%%%%%%%%%%%%%%%%%%%%%%%%%%
\section{Structure Generation Algorithm}
\label{sec:doctype-structure}
%%%%%%%%%%%%%%%%%%%%%%%%%%%%%%%%%%%%%%%%%%%%%%%%%%%%%%%%%%%%%%%%%%%%%%%%%%%%%%%

\textbf{Key Insight:} Document structure \textit{emerges} from the 8D coordinates.
Given a document type vector $\mathcal{D} = (D_1, ..., D_8)$, we can deterministically
generate the required structural elements (see Appendix UNMAPPED_DTE, Axioms DT-5 and DT-6).

\subsection{Structure Element Rules}

The following rules map dimensional thresholds to structural requirements:

\begin{table}[htbp]
\centering
\caption{Dimension-to-Structure Mapping Rules}
\label{tab:structure-rules}
\renewcommand{\arraystretch}{1.3}
\begin{tabular}{@{}lll@{}}
\toprule
\textbf{Condition} & \textbf{Structural Element} & \textbf{Rationale} \\
\midrule
$D_1 > 0.7$ & Technical Glossary & Experts need precise terminology \\
$D_1 < 0.3$ & Definition Box & Laypersons need explanations \\
$D_4 > 0.6$ & Detailed Subsections & Time permits depth \\
$D_4 < 0.3$ & Executive Summary & Time-constrained readers \\
$D_5 = G_1$ & Theory/Methods Section & Information requires foundation \\
$D_5 = G_2$ & Call-to-Action Box & Action needs clear next steps \\
$D_5 = G_3$ & Evidence Section & Persuasion needs proof \\
$D_7 > 0.7$ & Narrative Elements & Emotion needs story \\
$D_8 > 0.8$ & References + Version Control & Archival needs attribution \\
$D_8 > 0.8$ & Change Log & Archival needs history \\
$D_3 > 0.7$ & Stakeholder Analysis & Societal scope needs actors \\
$D_6 < 0.3$ & Confidentiality Notice & Internal needs protection \\
\bottomrule
\end{tabular}
\end{table}

\subsection{Computational Implementation}

\begin{verbatim}
from dataclasses import dataclass, field
from typing import List, Set
from enum import Enum

@dataclass
class StructuralElement:
    """A required or optional structural element."""
    name: str
    required: bool
    rationale: str
    template: str = ""

class DocumentStructure:
    """
    Emergent structure from 8D coordinates.

    Implements Axiom DT-5 (Structure Emergence) and
    DT-6 (Structural Element Rules) from Appendix UNMAPPED_DTE.
    """

    def __init__(self, doc_type: DocumentType):
        self.doc_type = doc_type
        self.elements: List[StructuralElement] = []
        self._generate_structure()

    def _generate_structure(self) -> None:
        """
        Apply dimension-to-structure mapping rules.

        Structure emerges deterministically from coordinates:
        S = h(D1, ..., D8)
        """
        d = self.doc_type

        # D1 Knowledge Rules
        if d.d1_knowledge > 0.7:
            self.elements.append(StructuralElement(
                name="Technical Glossary",
                required=True,
                rationale="Expert audience needs precise terminology",
                template="\\section{Glossary}\\n% Define technical terms"
            ))
        elif d.d1_knowledge < 0.3:
            self.elements.append(StructuralElement(
                name="Definition Box",
                required=True,
                rationale="Layperson audience needs concept explanations",
                template="\\begin{tcolorbox}[title=Key Terms]\\n\\end{tcolorbox}"
            ))

        # D4 Time Rules
        if d.d4_time > 0.6:
            self.elements.append(StructuralElement(
                name="Detailed Subsections",
                required=False,
                rationale="Audience has time for depth"
            ))
        if d.d4_time < 0.3:
            self.elements.append(StructuralElement(
                name="Executive Summary",
                required=True,
                rationale="Time-constrained readers need overview",
                template="\\begin{abstract}\\n% 150 words max\\n\\end{abstract}"
            ))

        # D5 Goal Rules
        if d.d5_goal == Goal.INFORM:
            self.elements.append(StructuralElement(
                name="Theory/Methods Section",
                required=True,
                rationale="Information transfer requires foundation",
                template="\\section{Theoretical Framework}\\n\\section{Methods}"
            ))
        elif d.d5_goal == Goal.ACT:
            self.elements.append(StructuralElement(
                name="Call-to-Action Box",
                required=True,
                rationale="Action-oriented needs clear next steps",
                template="\\begin{tcolorbox}[title=Next Steps]\\n\\end{tcolorbox}"
            ))
        elif d.d5_goal == Goal.PERSUADE:
            self.elements.append(StructuralElement(
                name="Evidence Section",
                required=True,
                rationale="Persuasion requires proof",
                template="\\section{Evidence}\\n% Data, studies, examples"
            ))

        # D7 Emotion Rules
        if d.d7_emotion > 0.7:
            self.elements.append(StructuralElement(
                name="Narrative Elements",
                required=True,
                rationale="High emotion benefits from storytelling"
            ))

        # D8 Persistence Rules
        if d.d8_persistence > 0.8:
            self.elements.append(StructuralElement(
                name="References Section",
                required=True,
                rationale="Archival documents need attribution",
                template="\\section{References}

% Link to master bibliography
\nocite{bcm_master}
\\n\\bibliographystyle{plain}"
            ))
            self.elements.append(StructuralElement(
                name="Version Control Header",
                required=True,
                rationale="Archival needs version tracking",
                template="% Version: X.Y (Date)\\n% Status: Draft/Final"
            ))

        # D3 Scope Rules
        if d.d3_scope > 0.7:
            self.elements.append(StructuralElement(
                name="Stakeholder Analysis",
                required=False,
                rationale="Societal scope involves multiple actors"
            ))

        # D6 Context Rules
        if d.d6_context < 0.3:
            self.elements.append(StructuralElement(
                name="Confidentiality Notice",
                required=True,
                rationale="Internal documents need protection",
                template="% CONFIDENTIAL - Internal Use Only"
            ))

    def get_required_elements(self) -> List[StructuralElement]:
        """Return only required structural elements."""
        return [e for e in self.elements if e.required]

    def get_optional_elements(self) -> List[StructuralElement]:
        """Return optional structural elements."""
        return [e for e in self.elements if not e.required]

    def generate_template(self) -> str:
        """Generate a document template with required structure."""
        template_parts = [
            f"% Document Type: {self.doc_type.name}",
            f"% Generated Structure from 8D: {self.doc_type.to_vector()}",
            "",
            "% === REQUIRED ELEMENTS ===",
        ]

        for elem in self.get_required_elements():
            template_parts.append(f"% {elem.name}: {elem.rationale}")
            if elem.template:
                template_parts.append(elem.template)
            template_parts.append("")

        if self.get_optional_elements():
            template_parts.append("% === OPTIONAL ELEMENTS ===")
            for elem in self.get_optional_elements():
                template_parts.append(f"% {elem.name}: {elem.rationale}")

        return "\n".join(template_parts)

    def validate_document(self, content: str) -> List[str]:
        """
        Validate that a document contains all required elements.

        Returns list of missing required elements.
        """
        missing = []
        for elem in self.get_required_elements():
            # Simple heuristic: check if element name appears
            if elem.name.lower() not in content.lower():
                missing.append(elem.name)
        return missing
\end{verbatim}

\subsection{Usage Example}

\begin{verbatim}
# Create a document type and generate its structure
policy_brief = DocumentType(
    name="Policy Brief",
    d1_knowledge=0.6,
    d2_proximity=0.4,
    d3_scope=0.85,
    d4_time=0.25,       # Low: time-constrained
    d5_goal=Goal.ACT,   # Action-oriented
    d6_context=1.0,
    d7_emotion=0.25,
    d8_persistence=0.7
)

# Structure emerges from coordinates
structure = DocumentStructure(policy_brief)

print("Required elements:")
for elem in structure.get_required_elements():
    print(f"  - {elem.name}: {elem.rationale}")

# Output:
# Required elements:
#   - Executive Summary: Time-constrained readers need overview
#   - Call-to-Action Box: Action-oriented needs clear next steps

# Generate template
template = structure.generate_template()
print(template)
\end{verbatim}

\begin{tcolorbox}[colback=yellow!5!white, colframe=yellow!75!black,
    title=Key Insight: Emergent Structure]

The structure generation algorithm implements the \textbf{audience-first principle}:

\begin{center}
\textbf{8D Coordinates} $\xrightarrow{h(\cdot)}$ \textbf{Document Structure}
\end{center}

Structure is not imposed top-down but emerges bottom-up from audience characteristics.
This explains why different document types have different structural conventions---
the structure serves the audience's needs as encoded in the 8D vector.

\end{tcolorbox}


%%%%%%%%%%%%%%%%%%%%%%%%%%%%%%%%%%%%%%%%%%%%%%%%%%%%%%%%%%%%%%%%%%%%%%%%%%%%%%%
\section{Style Generation Algorithm}
\label{sec:doctype-style}
%%%%%%%%%%%%%%%%%%%%%%%%%%%%%%%%%%%%%%%%%%%%%%%%%%%%%%%%%%%%%%%%%%%%%%%%%%%%%%%

\textbf{Key Insight:} Writing style \textit{emerges} from the 8D coordinates.
Given a document type vector $\mathcal{D} = (D_1, ..., D_8)$, we can deterministically
generate appropriate style parameters (see Appendix UNMAPPED_DTE, Axioms DT-7 and DT-8).

\subsection{Style Parameter Rules}

The following rules map dimensional thresholds to style requirements:

\begin{table}[htbp]
\centering
\caption{Dimension-to-Style Mapping Rules}
\label{tab:style-rules}
\renewcommand{\arraystretch}{1.3}
\begin{tabular}{@{}lll@{}}
\toprule
\textbf{Condition} & \textbf{Style Parameter} & \textbf{Value} \\
\midrule
$D_1 > 0.7$ & Flesch-Kincaid Grade & $\geq 14$ (graduate level) \\
$D_1 < 0.3$ & Flesch-Kincaid Grade & $\leq 8$ (middle school) \\
$D_2 > 0.8$ & Technical Jargon & Permitted \\
$D_4 > 0.6$ & Avg Sentence Length & $> 20$ words \\
$D_4 < 0.3$ & Avg Sentence Length & $< 15$ words \\
$D_6 > 0.7$ & Register & Formal \\
$D_6 < 0.3$ & Register & Informal \\
$D_7 > 0.6$ & Personal Pronouns & Yes (I/we) \\
$D_7 < 0.3$ & Personal Pronouns & No (impersonal) \\
$D_8 > 0.8$ & Hedging & Required (``suggests'') \\
$D_5 = G_3$ & Boosters & Permitted (``clearly'') \\
\bottomrule
\end{tabular}
\end{table}

\subsection{Computational Implementation}

\begin{verbatim}
from dataclasses import dataclass
from typing import List, Optional
from enum import Enum

class Register(Enum):
    INFORMAL = 1
    NEUTRAL = 2
    FORMAL = 3

@dataclass
class StyleParameter:
    """A style parameter with its value and rationale."""
    name: str
    value: any
    rationale: str

@dataclass
class StyleGuide:
    """Complete style guide generated from 8D coordinates."""
    flesch_kincaid_min: float
    flesch_kincaid_max: float
    avg_sentence_length_min: int
    avg_sentence_length_max: int
    register: Register
    personal_pronouns: bool
    hedging_required: bool
    boosters_permitted: bool
    technical_jargon: bool
    narrative_elements: bool
    active_voice_preferred: bool

class DocumentStyle:
    """
    Emergent style from 8D coordinates.

    Implements Axiom DT-7 (Style Emergence) and
    DT-8 (Style Parameter Rules) from Appendix UNMAPPED_DTE.
    """

    def __init__(self, doc_type: DocumentType):
        self.doc_type = doc_type
        self.parameters: List[StyleParameter] = []
        self.guide: StyleGuide = self._generate_style()

    def _generate_style(self) -> StyleGuide:
        """
        Apply dimension-to-style mapping rules.

        Style emerges deterministically from coordinates:
        sigma = s(D1, ..., D8)
        """
        d = self.doc_type

        # Default values
        fk_min, fk_max = 10, 14
        sent_min, sent_max = 15, 25
        register = Register.NEUTRAL
        pronouns = False
        hedging = False
        boosters = False
        jargon = False
        narrative = False
        active = True

        # D1 Knowledge Rules -> Vocabulary Complexity
        if d.d1_knowledge > 0.7:
            fk_min, fk_max = 14, 18
            self.parameters.append(StyleParameter(
                name="Flesch-Kincaid Grade",
                value="14-18 (graduate level)",
                rationale="Expert audience can handle complex text"
            ))
        elif d.d1_knowledge < 0.3:
            fk_min, fk_max = 6, 8
            self.parameters.append(StyleParameter(
                name="Flesch-Kincaid Grade",
                value="6-8 (middle school)",
                rationale="Layperson audience needs simple language"
            ))

        # D2 Proximity Rules -> Jargon
        if d.d2_proximity > 0.8:
            jargon = True
            self.parameters.append(StyleParameter(
                name="Technical Jargon",
                value="Permitted",
                rationale="Same-field audience shares vocabulary"
            ))
        elif d.d2_proximity < 0.3:
            jargon = False
            self.parameters.append(StyleParameter(
                name="Technical Jargon",
                value="Avoid",
                rationale="Cross-field audience needs translation"
            ))

        # D4 Time Rules -> Sentence Length
        if d.d4_time > 0.6:
            sent_min, sent_max = 20, 35
            self.parameters.append(StyleParameter(
                name="Sentence Length",
                value="20-35 words (complex)",
                rationale="Audience has time for nuanced prose"
            ))
        elif d.d4_time < 0.3:
            sent_min, sent_max = 8, 15
            self.parameters.append(StyleParameter(
                name="Sentence Length",
                value="8-15 words (short)",
                rationale="Time-constrained readers need brevity"
            ))

        # D6 Context Rules -> Register
        if d.d6_context > 0.7:
            register = Register.FORMAL
            active = False  # Passive acceptable
            self.parameters.append(StyleParameter(
                name="Register",
                value="Formal",
                rationale="External/public context requires formality"
            ))
        elif d.d6_context < 0.3:
            register = Register.INFORMAL
            active = True
            self.parameters.append(StyleParameter(
                name="Register",
                value="Informal",
                rationale="Internal context permits casual tone"
            ))

        # D7 Emotion Rules -> Pronouns & Narrative
        if d.d7_emotion > 0.6:
            pronouns = True
            narrative = True
            self.parameters.append(StyleParameter(
                name="Personal Pronouns",
                value="Yes (I/we)",
                rationale="Emotional engagement through personal voice"
            ))
        elif d.d7_emotion < 0.3:
            pronouns = False
            narrative = False
            self.parameters.append(StyleParameter(
                name="Personal Pronouns",
                value="No (impersonal)",
                rationale="Factual content requires objective voice"
            ))

        # D8 Persistence Rules -> Hedging
        if d.d8_persistence > 0.8:
            hedging = True
            self.parameters.append(StyleParameter(
                name="Hedging",
                value="Required",
                rationale="Archival documents need epistemic caution"
            ))

        # D5 Goal Rules -> Boosters
        if d.d5_goal == Goal.PERSUADE:
            boosters = True
            self.parameters.append(StyleParameter(
                name="Boosters",
                value="Permitted",
                rationale="Persuasion benefits from confidence markers"
            ))

        return StyleGuide(
            flesch_kincaid_min=fk_min,
            flesch_kincaid_max=fk_max,
            avg_sentence_length_min=sent_min,
            avg_sentence_length_max=sent_max,
            register=register,
            personal_pronouns=pronouns,
            hedging_required=hedging,
            boosters_permitted=boosters,
            technical_jargon=jargon,
            narrative_elements=narrative,
            active_voice_preferred=active
        )

    def generate_style_guide_text(self) -> str:
        """Generate human-readable style guide."""
        g = self.guide
        lines = [
            f"STYLE GUIDE: {self.doc_type.name}",
            f"Generated from 8D: {self.doc_type.to_vector()}",
            "",
            "VOCABULARY:",
            f"  - Flesch-Kincaid: {g.flesch_kincaid_min}-{g.flesch_kincaid_max}",
            f"  - Technical jargon: {'Permitted' if g.technical_jargon else 'Avoid'}",
            "",
            "SENTENCE STRUCTURE:",
            f"  - Length: {g.avg_sentence_length_min}-{g.avg_sentence_length_max} words",
            f"  - Voice: {'Active preferred' if g.active_voice_preferred else 'Passive acceptable'}",
            "",
            "TONE:",
            f"  - Register: {g.register.name}",
            f"  - Personal pronouns: {'Yes' if g.personal_pronouns else 'No'}",
            f"  - Narrative elements: {'Yes' if g.narrative_elements else 'No'}",
            "",
            "EPISTEMIC MARKERS:",
            f"  - Hedging: {'Required' if g.hedging_required else 'Optional'}",
            f"  - Boosters: {'Permitted' if g.boosters_permitted else 'Avoid'}",
        ]
        return "\n".join(lines)

    def validate_text(self, text: str) -> List[str]:
        """
        Validate text against style guide.

        Returns list of style violations.
        """
        import re
        violations = []

        # Check sentence length (approximate)
        sentences = re.split(r'[.!?]+', text)
        sentences = [s.strip() for s in sentences if s.strip()]
        if sentences:
            avg_len = sum(len(s.split()) for s in sentences) / len(sentences)
            if avg_len < self.guide.avg_sentence_length_min:
                violations.append(
                    f"Avg sentence length {avg_len:.1f} < min {self.guide.avg_sentence_length_min}"
                )
            if avg_len > self.guide.avg_sentence_length_max:
                violations.append(
                    f"Avg sentence length {avg_len:.1f} > max {self.guide.avg_sentence_length_max}"
                )

        # Check personal pronouns
        pronoun_pattern = r'\b(I|we|my|our)\b'
        has_pronouns = bool(re.search(pronoun_pattern, text, re.IGNORECASE))
        if has_pronouns and not self.guide.personal_pronouns:
            violations.append("Personal pronouns used but not permitted")
        if not has_pronouns and self.guide.personal_pronouns:
            violations.append("No personal pronouns found but recommended")

        return violations
\end{verbatim}

\subsection{Usage Example}

\begin{verbatim}
# Create document type and generate style
bcm_appendix = DocumentType(
    name="EBF Appendix",
    d1_knowledge=0.85,
    d2_proximity=0.75,
    d3_scope=0.6,
    d4_time=0.7,
    d5_goal=Goal.INFORM,
    d6_context=1.0,
    d7_emotion=0.15,
    d8_persistence=0.9
)

# Style emerges from coordinates
style = DocumentStyle(bcm_appendix)

# Print style guide
print(style.generate_style_guide_text())

# Output:
# STYLE GUIDE: EBF Appendix
# Generated from 8D: [0.85, 0.75, 0.6, 0.7, 0.14, 1.0, 0.15, 0.9]
#
# VOCABULARY:
#   - Flesch-Kincaid: 14-18
#   - Technical jargon: Permitted
#
# SENTENCE STRUCTURE:
#   - Length: 20-35 words
#   - Voice: Passive acceptable
#
# TONE:
#   - Register: FORMAL
#   - Personal pronouns: No
#   - Narrative elements: No
#
# EPISTEMIC MARKERS:
#   - Hedging: Required
#   - Boosters: Avoid

# Validate a text sample
violations = style.validate_text(sample_text)
\end{verbatim}

\begin{tcolorbox}[colback=yellow!5!white, colframe=yellow!75!black,
    title=Key Insight: Emergent Style]

The style generation algorithm extends the \textbf{emergence chain}:

\begin{center}
\textbf{8D Coordinates} $\xrightarrow{h(\cdot)}$ \textbf{Structure} $\xrightarrow{s(\cdot)}$ \textbf{Style}
\end{center}

Both structure and style are not arbitrary conventions but functional responses
to audience needs encoded in the 8D vector. This explains why:
\begin{itemize}[nosep]
\item Academic papers use hedging (high $D_8$)
\item Tweets are short (low $D_4$)
\item Love letters use ``I'' (high $D_7$)
\end{itemize}

\end{tcolorbox}


%%%%%%%%%%%%%%%%%%%%%%%%%%%%%%%%%%%%%%%%%%%%%%%%%%%%%%%%%%%%%%%%%%%%%%%%%%%%%%%
\section{Vocabulary Generation Algorithm}
\label{sec:doctype-vocabulary}
%%%%%%%%%%%%%%%%%%%%%%%%%%%%%%%%%%%%%%%%%%%%%%%%%%%%%%%%%%%%%%%%%%%%%%%%%%%%%%%

\textbf{Key Insight:} Vocabulary \textit{emerges} from style parameters.
Given a style guide $\sigma$, we can deterministically generate appropriate
word class constraints (see Appendix UNMAPPED_DTE, Axiom DT-9).

\subsection{Vocabulary Class Rules}

The following rules map style parameters to vocabulary constraints:

\begin{table}[htbp]
\centering
\caption{Style-to-Vocabulary Mapping Rules}
\label{tab:vocabulary-rules}
\renewcommand{\arraystretch}{1.3}
\begin{tabular}{@{}lll@{}}
\toprule
\textbf{Style Parameter} & \textbf{Word Class} & \textbf{Examples} \\
\midrule
$\sigma_{\text{hedging}}=\text{True}$ & Hedging words & suggests, may, indicates, appears \\
$\sigma_{\text{hedging}}=\text{False}$ & Assertive words & shows, proves, demonstrates, is \\
$\sigma_{\text{boosters}}=\text{True}$ & Booster words & clearly, certainly, undoubtedly \\
$\sigma_{\text{register}}=\text{Formal}$ & Formal connectives & therefore, consequently, however \\
$\sigma_{\text{register}}=\text{Informal}$ & Informal connectives & so, but, anyway \\
$\sigma_{\text{jargon}}=\text{True}$ & Technical terms & Field-specific vocabulary \\
$\sigma_{\text{pronouns}}=\text{True}$ & Personal pronouns & I, we, my, our \\
$\sigma_{\text{pronouns}}=\text{False}$ & Impersonal references & the study, results, analysis \\
\bottomrule
\end{tabular}
\end{table}

\subsection{Computational Implementation}

\begin{verbatim}
from dataclasses import dataclass, field
from typing import List, Set, Dict
from enum import Enum

@dataclass
class WordClass:
    """A vocabulary constraint with word lists."""
    name: str
    words: Set[str]
    required: bool  # Must use these words
    forbidden: bool  # Must NOT use these words
    rationale: str

@dataclass
class VocabularyGuide:
    """Complete vocabulary guide generated from style parameters."""
    hedging_words: Set[str] = field(default_factory=set)
    assertive_words: Set[str] = field(default_factory=set)
    booster_words: Set[str] = field(default_factory=set)
    formal_connectives: Set[str] = field(default_factory=set)
    informal_connectives: Set[str] = field(default_factory=set)
    personal_pronouns: Set[str] = field(default_factory=set)
    impersonal_references: Set[str] = field(default_factory=set)
    technical_terms_permitted: bool = False

    # Constraints
    use_hedging: bool = False
    use_boosters: bool = False
    use_formal_connectives: bool = True
    use_personal_pronouns: bool = False

class DocumentVocabulary:
    """
    Emergent vocabulary from style parameters.

    Implements Axiom DT-9 (Vocabulary Emergence) from Appendix UNMAPPED_DTE.
    """

    # Standard word lists
    HEDGING_WORDS = {
        "suggests", "may", "might", "could", "appears", "seems",
        "indicates", "potentially", "possibly", "likely", "tends to"
    }
    ASSERTIVE_WORDS = {
        "shows", "proves", "demonstrates", "is", "reveals",
        "confirms", "establishes", "determines"
    }
    BOOSTER_WORDS = {
        "clearly", "certainly", "undoubtedly", "definitely",
        "obviously", "evidently", "surely", "absolutely"
    }
    FORMAL_CONNECTIVES = {
        "therefore", "consequently", "however", "furthermore",
        "moreover", "nevertheless", "thus", "hence", "accordingly"
    }
    INFORMAL_CONNECTIVES = {
        "so", "but", "and", "also", "anyway", "besides",
        "still", "though", "yet"
    }
    PERSONAL_PRONOUNS = {"I", "we", "my", "our", "me", "us"}
    IMPERSONAL_REFS = {
        "the study", "the analysis", "results", "findings",
        "the data", "the evidence", "this paper", "the research"
    }

    def __init__(self, style: DocumentStyle):
        self.style = style
        self.word_classes: List[WordClass] = []
        self.guide: VocabularyGuide = self._generate_vocabulary()

    def _generate_vocabulary(self) -> VocabularyGuide:
        """
        Apply style-to-vocabulary mapping rules.

        Vocabulary emerges from style parameters:
        V = v(sigma)
        """
        g = self.style.guide
        guide = VocabularyGuide()

        # Hedging rules
        if g.hedging_required:
            guide.hedging_words = self.HEDGING_WORDS
            guide.use_hedging = True
            self.word_classes.append(WordClass(
                name="Hedging Words",
                words=self.HEDGING_WORDS,
                required=True,
                forbidden=False,
                rationale="Archival documents need epistemic caution"
            ))
        else:
            guide.assertive_words = self.ASSERTIVE_WORDS
            guide.use_hedging = False
            self.word_classes.append(WordClass(
                name="Assertive Words",
                words=self.ASSERTIVE_WORDS,
                required=False,
                forbidden=False,
                rationale="Non-archival allows direct assertions"
            ))

        # Booster rules
        if g.boosters_permitted:
            guide.booster_words = self.BOOSTER_WORDS
            guide.use_boosters = True
            self.word_classes.append(WordClass(
                name="Booster Words",
                words=self.BOOSTER_WORDS,
                required=False,
                forbidden=False,
                rationale="Persuasion benefits from confidence markers"
            ))
        else:
            self.word_classes.append(WordClass(
                name="Booster Words",
                words=self.BOOSTER_WORDS,
                required=False,
                forbidden=True,
                rationale="Avoid overconfidence in non-persuasive text"
            ))

        # Register rules
        if g.register == Register.FORMAL:
            guide.formal_connectives = self.FORMAL_CONNECTIVES
            guide.use_formal_connectives = True
            self.word_classes.append(WordClass(
                name="Formal Connectives",
                words=self.FORMAL_CONNECTIVES,
                required=True,
                forbidden=False,
                rationale="Formal register requires academic connectives"
            ))
        elif g.register == Register.INFORMAL:
            guide.informal_connectives = self.INFORMAL_CONNECTIVES
            guide.use_formal_connectives = False
            self.word_classes.append(WordClass(
                name="Informal Connectives",
                words=self.INFORMAL_CONNECTIVES,
                required=False,
                forbidden=False,
                rationale="Informal register permits casual connectives"
            ))

        # Pronoun rules
        if g.personal_pronouns:
            guide.personal_pronouns = self.PERSONAL_PRONOUNS
            guide.use_personal_pronouns = True
            self.word_classes.append(WordClass(
                name="Personal Pronouns",
                words=self.PERSONAL_PRONOUNS,
                required=True,
                forbidden=False,
                rationale="Emotional engagement through personal voice"
            ))
        else:
            guide.impersonal_references = self.IMPERSONAL_REFS
            guide.use_personal_pronouns = False
            self.word_classes.append(WordClass(
                name="Impersonal References",
                words=self.IMPERSONAL_REFS,
                required=True,
                forbidden=False,
                rationale="Objective voice requires impersonal style"
            ))

        # Jargon rules
        guide.technical_terms_permitted = g.technical_jargon

        return guide

    def generate_vocabulary_guide_text(self) -> str:
        """Generate human-readable vocabulary guide."""
        g = self.guide
        lines = [
            f"VOCABULARY GUIDE: {self.style.doc_type.name}",
            f"Generated from Style: {self.style.guide}",
            "",
            "EPISTEMIC MARKERS:",
        ]

        if g.use_hedging:
            lines.append(f"  - USE hedging: {', '.join(sorted(g.hedging_words)[:5])}...")
        else:
            lines.append(f"  - Assertive OK: {', '.join(sorted(g.assertive_words)[:5])}...")

        if g.use_boosters:
            lines.append(f"  - Boosters OK: {', '.join(sorted(g.booster_words)[:5])}...")
        else:
            lines.append("  - AVOID boosters")

        lines.extend([
            "",
            "CONNECTIVES:",
        ])
        if g.use_formal_connectives:
            lines.append(f"  - USE formal: {', '.join(sorted(g.formal_connectives)[:5])}...")
        else:
            lines.append(f"  - Informal OK: {', '.join(sorted(g.informal_connectives)[:5])}...")

        lines.extend([
            "",
            "SUBJECT REFERENCES:",
        ])
        if g.use_personal_pronouns:
            lines.append(f"  - USE personal: {', '.join(sorted(g.personal_pronouns))}")
        else:
            lines.append(f"  - USE impersonal: {', '.join(sorted(g.impersonal_references)[:4])}...")

        lines.extend([
            "",
            "TECHNICAL TERMS:",
            f"  - Field jargon: {'Permitted' if g.technical_terms_permitted else 'Avoid'}",
        ])

        return "\n".join(lines)

    def validate_text(self, text: str) -> List[str]:
        """
        Validate text against vocabulary constraints.

        Returns list of vocabulary violations.
        """
        import re
        violations = []
        text_lower = text.lower()

        # Check hedging requirement
        if self.guide.use_hedging:
            hedge_found = any(w in text_lower for w in self.guide.hedging_words)
            if not hedge_found:
                violations.append("No hedging words found (required for archival)")

        # Check booster restriction
        if not self.guide.use_boosters:
            boosters_found = [w for w in self.BOOSTER_WORDS if w in text_lower]
            if boosters_found:
                violations.append(f"Boosters used but not permitted: {boosters_found}")

        # Check pronoun usage
        pronoun_pattern = r'\b(I|we|my|our)\b'
        has_pronouns = bool(re.search(pronoun_pattern, text, re.IGNORECASE))
        if has_pronouns and not self.guide.use_personal_pronouns:
            violations.append("Personal pronouns used but impersonal style required")
        if not has_pronouns and self.guide.use_personal_pronouns:
            violations.append("No personal pronouns found but personal style required")

        return violations
\end{verbatim}

\subsection{Usage Example}

\begin{verbatim}
# Create document type -> style -> vocabulary
bcm_appendix = DocumentType(
    name="EBF Appendix",
    d1_knowledge=0.85,
    d2_proximity=0.75,
    d3_scope=0.6,
    d4_time=0.7,
    d5_goal=Goal.INFORM,
    d6_context=1.0,
    d7_emotion=0.15,
    d8_persistence=0.9
)

# Chain: 8D -> Style -> Vocabulary
style = DocumentStyle(bcm_appendix)
vocabulary = DocumentVocabulary(style)

# Print vocabulary guide
print(vocabulary.generate_vocabulary_guide_text())

# Output:
# VOCABULARY GUIDE: EBF Appendix
# Generated from Style: StyleGuide(...)
#
# EPISTEMIC MARKERS:
#   - USE hedging: appears, could, indicates, likely, may...
#   - AVOID boosters
#
# CONNECTIVES:
#   - USE formal: accordingly, consequently, furthermore, hence, however...
#
# SUBJECT REFERENCES:
#   - USE impersonal: findings, results, the analysis, the data...
#
# TECHNICAL TERMS:
#   - Field jargon: Permitted

# Validate a text sample
violations = vocabulary.validate_text(sample_text)
\end{verbatim}

\begin{tcolorbox}[colback=yellow!5!white, colframe=yellow!75!black,
    title=Key Insight: Complete Emergence Chain]

The output generation algorithm completes the \textbf{full emergence chain}:

\begin{center}
\textbf{8D Coordinates} $\xrightarrow{h(\cdot)}$ \textbf{Structure} $\xrightarrow{s(\cdot)}$ \textbf{Style} $\xrightarrow{v(\cdot)}$ \textbf{Vocabulary} $\xrightarrow{o(\cdot)}$ \textbf{Output (PDF)}
\end{center}

This implements the \textbf{audience-first principle} completely:
\begin{enumerate}[nosep]
\item Specify your audience $\to$ derive 8D coordinates
\item 8D $\to$ document structure (what sections?)
\item 8D $\to$ style parameters (how to write?)
\item Style $\to$ vocabulary constraints (what words?)
\item \textbf{Structure + Style + Vocabulary $\to$ Output format + compiled document}
\end{enumerate}

\textbf{Practical implication:} Given ``who is my reader?'', the system generates
the complete document including compiled PDF output.

\end{tcolorbox}


%%%%%%%%%%%%%%%%%%%%%%%%%%%%%%%%%%%%%%%%%%%%%%%%%%%%%%%%%%%%%%%%%%%%%%%%%%%%%%%
\section{Output Generation Algorithm}
\label{sec:doctype-output}
%%%%%%%%%%%%%%%%%%%%%%%%%%%%%%%%%%%%%%%%%%%%%%%%%%%%%%%%%%%%%%%%%%%%%%%%%%%%%%%

The final step in the emergence chain transforms the generated structure, style,
and vocabulary into a \textbf{compiled output document}. This section specifies
the output format defaults and automatic compilation process.

\subsection{Output Format Rules}

Output format is determined by document type characteristics:

\begin{table}[htbp]
\centering
\caption{Output Format Defaults by Document Type}
\label{tab:output-format-defaults}
\renewcommand{\arraystretch}{1.3}
\begin{tabular}{@{}llllc@{}}
\toprule
\textbf{Document Type} & \textbf{$D_8$} & \textbf{Format} & \textbf{Template} & \textbf{Auto-PDF} \\
\midrule
Journal Paper & $> 0.9$ & LaTeX & AER/JPE/QJE & \checkmark \\
Working Paper & $> 0.8$ & LaTeX & NBER/SSRN & \checkmark \\
EBF Appendix & $> 0.85$ & LaTeX & EBF-Template & \checkmark \\
Policy Brief & $> 0.6$ & LaTeX/MD & Policy-Template & \checkmark \\
Technical Report & $> 0.7$ & LaTeX & Report-Template & \checkmark \\
Internal Memo & $< 0.5$ & Markdown & Memo-Template & -- \\
Blog Post & $< 0.4$ & Markdown & Blog-Template & -- \\
Email & $< 0.3$ & Plain Text & -- & -- \\
\bottomrule
\end{tabular}
\end{table}

\textbf{Rule UNMAPPED_AUT-1 (Format Selection):}
\begin{equation}
\text{Format}(D_8, D_6) = \begin{cases}
\text{LaTeX} & \text{if } D_8 > 0.6 \text{ or } (D_6 > 0.8 \text{ and } D_1 > 0.7) \\
\text{Markdown} & \text{if } 0.3 < D_8 \leq 0.6 \\
\text{Plain Text} & \text{if } D_8 \leq 0.3
\end{cases}
\end{equation}

\textbf{Rule UNMAPPED_AUT-2 (Auto-Compile):}
\begin{equation}
\text{AutoCompile} = \begin{cases}
\text{True} & \text{if Format} = \text{LaTeX and } D_8 > 0.5 \\
\text{False} & \text{otherwise}
\end{cases}
\end{equation}

\subsection{Template Selection}

Templates are selected based on document purpose ($D_5$) and audience ($D_1$, $D_2$):

\begin{table}[htbp]
\centering
\caption{Template Selection by Purpose and Audience}
\label{tab:template-selection}
\renewcommand{\arraystretch}{1.3}
\begin{tabular}{@{}llll@{}}
\toprule
\textbf{$D_5$ (Goal)} & \textbf{$D_1 > 0.7$} & \textbf{$D_1 \leq 0.7$} & \textbf{Style Variant} \\
\midrule
$G_1$ (Inform) & Academic & Popular Science & Fehr, Thaler \\
$G_2$ (Action) & Policy Brief & Call-to-Action & Sunstein \\
$G_3$ (Persuade) & Review Article & Op-Ed & Kahneman \\
\bottomrule
\end{tabular}
\end{table}

\subsection{Computational Implementation}

\begin{verbatim}
from dataclasses import dataclass
from enum import Enum
from typing import Optional
import subprocess
import os

class OutputFormat(Enum):
    LATEX = "latex"
    MARKDOWN = "markdown"
    PLAIN = "plain"

class Template(Enum):
    AER = "aer"           # American Economic Review style
    JPE = "jpe"           # Journal of Political Economy
    NBER = "nber"         # NBER Working Paper
    EBF = "ebf"           # EBF Framework style
    POLICY = "policy"     # Policy Brief
    FEHR = "fehr"         # Ernst Fehr style (Zurich)
    THALER = "thaler"     # Richard Thaler style (Chicago)

@dataclass
class OutputConfig:
    """Output configuration derived from 8D coordinates."""
    format: OutputFormat
    template: Template
    auto_compile: bool
    output_dir: str = "outputs"

    @classmethod
    def from_8d(cls, d1: float, d5: str, d6: float, d8: float,
                style_variant: str = "fehr") -> "OutputConfig":
        """Derive output configuration from 8D coordinates."""

        # Rule UNMAPPED_AUT-1: Format Selection
        if d8 > 0.6 or (d6 > 0.8 and d1 > 0.7):
            fmt = OutputFormat.LATEX
        elif d8 > 0.3:
            fmt = OutputFormat.MARKDOWN
        else:
            fmt = OutputFormat.PLAIN

        # Rule UNMAPPED_AUT-2: Auto-Compile
        auto_compile = (fmt == OutputFormat.LATEX and d8 > 0.5)

        # Template Selection
        if style_variant.lower() == "fehr":
            template = Template.FEHR
        elif d5 == "G1" and d1 > 0.7:
            template = Template.AER
        elif d5 == "G2":
            template = Template.POLICY
        else:
            template = Template.EBF

        return cls(
            format=fmt,
            template=template,
            auto_compile=auto_compile
        )

@dataclass
class OutputGenerator:
    """Generates and compiles document output."""
    config: OutputConfig

    def generate(self, content: str, filename: str) -> str:
        """Generate output file and optionally compile to PDF."""

        # Ensure output directory exists
        os.makedirs(self.config.output_dir, exist_ok=True)

        if self.config.format == OutputFormat.LATEX:
            return self._generate_latex(content, filename)
        elif self.config.format == OutputFormat.MARKDOWN:
            return self._generate_markdown(content, filename)
        else:
            return self._generate_plain(content, filename)

    def _generate_latex(self, content: str, filename: str) -> str:
        """Generate LaTeX file and compile to PDF if auto_compile."""

        tex_path = f"{self.config.output_dir}/{filename}.tex"
        pdf_path = f"{self.config.output_dir}/{filename}.pdf"

        # Write LaTeX file
        with open(tex_path, 'w') as f:
            f.write(content)

        # Auto-compile if configured
        if self.config.auto_compile:
            self._compile_pdf(tex_path)
            return pdf_path

        return tex_path

    def _compile_pdf(self, tex_path: str) -> bool:
        """Compile LaTeX to PDF using pdflatex."""
        try:
            # Run pdflatex twice for references
            for _ in range(2):
                subprocess.run(
                    ["pdflatex", "-interaction=nonstopmode", tex_path],
                    cwd=self.config.output_dir,
                    capture_output=True,
                    check=True
                )
            return True
        except subprocess.CalledProcessError:
            return False

    def _generate_markdown(self, content: str, filename: str) -> str:
        """Generate Markdown file."""
        md_path = f"{self.config.output_dir}/{filename}.md"
        with open(md_path, 'w') as f:
            f.write(content)
        return md_path

    def _generate_plain(self, content: str, filename: str) -> str:
        """Generate plain text file."""
        txt_path = f"{self.config.output_dir}/{filename}.txt"
        with open(txt_path, 'w') as f:
            f.write(content)
        return txt_path
\end{verbatim}

\subsection{Integration with Document Pipeline}

The complete document generation pipeline:

\begin{verbatim}
def generate_document(audience_description: str,
                      topic: str,
                      style_variant: str = "fehr") -> str:
    """
    Complete 8D document generation pipeline.

    8D -> Structure -> Style -> Vocabulary -> Output (PDF)
    """

    # Step 1: Measure 8D coordinates from audience description
    coords = measure_8d_with_llm(audience_description)

    # Step 2: Generate structure
    structure = DocumentStructure.from_8d(
        d1=coords.d1, d4=coords.d4, d5=coords.d5,
        d6=coords.d6, d8=coords.d8
    )
    template = structure.generate_template()

    # Step 3: Generate style
    style = DocumentStyle.from_8d(
        d1=coords.d1, d4=coords.d4, d6=coords.d6,
        d7=coords.d7, d8=coords.d8
    )

    # Step 4: Generate vocabulary
    vocab = VocabularyProfile.from_style(style)

    # Step 5: Configure output
    output_config = OutputConfig.from_8d(
        d1=coords.d1, d5=coords.d5, d6=coords.d6, d8=coords.d8,
        style_variant=style_variant
    )

    # Step 6: Generate content (LLM call with constraints)
    content = generate_content_with_llm(
        topic=topic,
        template=template,
        style=style,
        vocab=vocab
    )

    # Step 7: Output and compile
    generator = OutputGenerator(output_config)
    output_path = generator.generate(content, filename=topic.replace(" ", "_"))

    return output_path

# Example usage
pdf_path = generate_document(
    audience_description="Behavioral economists at top research universities",
    topic="Behavioral Change Journey",
    style_variant="fehr"
)
print(f"Generated: {pdf_path}")
# Output: Generated: outputs/Behavioral_Change_Journey.pdf
\end{verbatim}

\subsection{Style Variants}

The system supports named style variants based on prominent authors:

\begin{table}[htbp]
\centering
\caption{Named Style Variants}
\label{tab:style-variants}
\renewcommand{\arraystretch}{1.3}
\begin{tabular}{@{}llp{6cm}@{}}
\toprule
\textbf{Variant} & \textbf{Template} & \textbf{Characteristics} \\
\midrule
\texttt{fehr} & AER/JPE & Formal models, testable hypotheses, experimental design, complementarity focus \\
\texttt{thaler} & AER/JPE & Accessible, anomalies-first, practical implications \\
\texttt{kahneman} & Science/Nature & Dual-system framing, heuristics, broad appeal \\
\texttt{sunstein} & Policy & Regulatory focus, cost-benefit, implementation \\
\bottomrule
\end{tabular}
\end{table}

\begin{tcolorbox}[colback=green!5!white, colframe=green!75!black,
    title=Axiom UNMAPPED_AUT-1: Output Completeness]
A document generation is \textbf{complete} if and only if:
\begin{enumerate}[nosep]
\item 8D coordinates are specified
\item Structure, Style, and Vocabulary are derived
\item Output format is determined by Rules UNMAPPED_AUT-1, UNMAPPED_AUT-2
\item For LaTeX documents with $D_8 > 0.5$: PDF is automatically compiled
\end{enumerate}
\end{tcolorbox}


%%%%%%%%%%%%%%%%%%%%%%%%%%%%%%%%%%%%%%%%%%%%%%%%%%%%%%%%%%%%%%%%%%%%%%%%%%%%%%%
\section{Worked Example}
\label{sec:doctype-example}
%%%%%%%%%%%%%%%%%%%%%%%%%%%%%%%%%%%%%%%%%%%%%%%%%%%%%%%%%%%%%%%%%%%%%%%%%%%%%%%

\subsection{Example: Characterizing a EBF Appendix}

\paragraph{Setup.} We want to characterize the target document type for a
EBF appendix (like this one).

\paragraph{Application.}

\begin{verbatim}
bcm_appendix = DocumentType(
    name="EBF Appendix",
    d1_knowledge=0.85,    # Expert: Behavioral economists, researchers
    d2_proximity=0.75,    # Close: Economics, psychology, policy
    d3_scope=0.6,         # Field-wide, some societal
    d4_time=0.7,          # Detailed reading expected
    d5_goal=Goal.INFORM,  # Primary: knowledge transfer
    d6_context=1.0,       # Public/academic
    d7_emotion=0.15,      # Mostly factual, some engagement
    d8_persistence=0.9    # Archival, reference material
)

# Find similar document types
nearest = find_nearest_type(bcm_appendix, DOCUMENT_TYPES, n=3)
print("Nearest types:", nearest)
# Output: [('Journal Paper', 0.21), ('Policy Brief', 0.45), ...]
\end{verbatim}

\paragraph{Result.} The EBF appendix is closest to ``Journal Paper'' but
with slightly lower knowledge requirements (0.85 vs 0.9) and higher practical
orientation.

\paragraph{Discussion.} This analysis suggests that EBF appendices should:
\begin{itemize}[nosep]
\item Use formal academic style (high $D_1$, $D_2$)
\item Include practical examples (moderate $D_3$)
\item Be thorough but accessible (moderate $D_4$)
\item Maintain factual tone with engagement (low but non-zero $D_7$)
\end{itemize}


%%%%%%%%%%%%%%%%%%%%%%%%%%%%%%%%%%%%%%%%%%%%%%%%%%%%%%%%%%%%%%%%%%%%%%%%%%%%%%%
\section{Emergent Clusters}
\label{sec:doctype-clusters}
%%%%%%%%%%%%%%%%%%%%%%%%%%%%%%%%%%%%%%%%%%%%%%%%%%%%%%%%%%%%%%%%%%%%%%%%%%%%%%%

When clustering a large corpus of document types, natural groupings emerge:

\begin{table}[htbp]
\centering
\caption{Emergent Document Clusters}
\label{tab:clusters}
\renewcommand{\arraystretch}{1.3}
\begin{tabular}{@{}lp{6cm}l@{}}
\toprule
\textbf{Cluster} & \textbf{Characteristics} & \textbf{Examples} \\
\midrule
\textbf{Scientific} &
High $D_1$, $D_2$; $D_5$=Inform; Low $D_7$; High $D_8$ &
Journal, Thesis, Review \\
\textbf{Policy} &
Medium $D_1$; $D_5$=Act; High $D_3$; Low $D_4$ &
Brief, Report, Memo \\
\textbf{Journalistic} &
Low-Medium $D_1$; $D_5$=Inform; Low $D_4$; Low $D_8$ &
News, Feature, Blog \\
\textbf{Creative} &
Variable $D_1$; $D_5$=Entertain/Experience; High $D_7$ &
Novel, Poetry, Film \\
\textbf{Personal} &
Low $D_3$, $D_6$; $D_5$=Connect/Express; High $D_7$ &
Letter, Diary, Message \\
\textbf{Technical} &
High $D_1$, $D_2$; $D_5$=Act; Low $D_7$; High $D_8$ &
Manual, Spec, API Doc \\
\bottomrule
\end{tabular}
\end{table}

\textbf{Key Insight:} These clusters are \textit{emergent}---they arise from
the data, not from predefined categories. The names are post-hoc labels for
regions in the 8D space that happen to be densely populated.


%%%%%%%%%%%%%%%%%%%%%%%%%%%%%%%%%%%%%%%%%%%%%%%%%%%%%%%%%%%%%%%%%%%%%%%%%%%%%%%
\section{Implications for Practice}
\label{sec:doctype-practice}
%%%%%%%%%%%%%%%%%%%%%%%%%%%%%%%%%%%%%%%%%%%%%%%%%%%%%%%%%%%%%%%%%%%%%%%%%%%%%%%

\begin{tcolorbox}[colback=green!5!white, colframe=green!75!black,
    title=Practical Implications]
\begin{enumerate}
\item \textbf{Document Design:} Before writing, specify the 8D target vector.
      This determines appropriate style, length, and tone.

\item \textbf{Style Validation:} Use LLM measurement to verify that a document
      matches its intended type. Flag deviations for revision.

\item \textbf{Gap Identification:} Analyze which regions of the 8D space lack
      document types. These represent communication opportunities.

\item \textbf{Audience Analysis:} Instead of asking ``Who is my audience?''
      ask ``What are my audience's 8D coordinates?''

\item \textbf{Cross-Cultural Adaptation:} Same content, different $D_1$/$D_4$/$D_7$
      settings for different audiences.
\end{enumerate}
\end{tcolorbox}


%%%%%%%%%%%%%%%%%%%%%%%%%%%%%%%%%%%%%%%%%%%%%%%%%%%%%%%%%%%%%%%%%%%%%%%%%%%%%%%
\section{Integration with Other Appendices}
\label{sec:doctype-integration}
%%%%%%%%%%%%%%%%%%%%%%%%%%%%%%%%%%%%%%%%%%%%%%%%%%%%%%%%%%%%%%%%%%%%%%%%%%%%%%%

\subsection{Connection to CORE-WHEN (V): Context Parallels}

\begin{tcolorbox}[colback=yellow!5!white, colframe=yellow!75!black,
    title=Integration: METHOD-DOCTYPE $\leftrightarrow$ CORE-WHEN]

\textbf{Parallel Structure:}

The 8D document dimensions parallel the $\Psi$ context dimensions:

\begin{center}
\begin{tabular}{ll}
$D_3$ (Scope) & $\leftrightarrow$ $\Psi_{Social}$ \\
$D_4$ (Time) & $\leftrightarrow$ $\Psi_{Temporal}$ \\
$D_6$ (Context) & $\leftrightarrow$ $\Psi_{Institutional}$ \\
$D_7$ (Emotion) & $\leftrightarrow$ $\Psi_{Cultural}$ \\
\end{tabular}
\end{center}

\textbf{Implication:} Document type selection should consider how $\Psi$ affects
the audience's reception of different document styles.

\end{tcolorbox}

\subsection{Connection to METHOD-LLMMC (AN)}

The LLM-based measurement approach in Section~\ref{sec:doctype-llm} directly
applies the Monte Carlo methodology from Appendix LLM1. Key adaptations:

\begin{itemize}[nosep]
\item \textbf{Prompt Engineering}: Domain-specific prompts for document analysis
\item \textbf{Aggregation}: Median for continuous, mode for categorical
\item \textbf{Validation}: Compare LLM ratings to expert ratings for calibration
\end{itemize}


%%%%%%%%%%%%%%%%%%%%%%%%%%%%%%%%%%%%%%%%%%%%%%%%%%%%%%%%%%%%%%%%%%%%%%%%%%%%%%%
\section{Summary}
\label{sec:doctype-summary}
%%%%%%%%%%%%%%%%%%%%%%%%%%%%%%%%%%%%%%%%%%%%%%%%%%%%%%%%%%%%%%%%%%%%%%%%%%%%%%%

\begin{tcolorbox}[colback=green!10!white, colframe=green!75!black,
    title=\textbf{Appendix UNMAPPED_DTC: Summary}]

\textbf{Core Question:} How to systematically characterize document types?

\textbf{Main Contribution:}
\begin{itemize}[nosep]
    \item 8D audience-centered document space
    \item Computational clustering framework
    \item LLM-based dimension measurement
    \item \textbf{Complete emergence chain}: Structure $\to$ Style $\to$ Vocabulary
\end{itemize}

\textbf{Key Outputs:}
\begin{itemize}[nosep]
    \item $\mathcal{D} = (D_1, ..., D_8)$: Document type vector
    \item Clustering algorithms for emergent type discovery
    \item Gap analysis for communication opportunities
    \item $\mathcal{S} = h(D_1, ..., D_8)$: Emergent document structure
    \item $\sigma = s(D_1, ..., D_8)$: Emergent writing style
    \item $\mathcal{V} = v(\sigma)$: Emergent vocabulary constraints
\end{itemize}

\textbf{Integration:} Connects to CORE-WHEN ($\Psi$ parallels),
METHOD-LLMMC (measurement methodology), and REF-DOCTYPE (DDD, axioms DT-1 to DT-9).

\end{tcolorbox}


%%%%%%%%%%%%%%%%%%%%%%%%%%%%%%%%%%%%%%%%%%%%%%%%%%%%%%%%%%%%%%%%%%%%%%%%%%%%%%%
\section{Glossary of Symbols}
\label{sec:doctype-glossary}
%%%%%%%%%%%%%%%%%%%%%%%%%%%%%%%%%%%%%%%%%%%%%%%%%%%%%%%%%%%%%%%%%%%%%%%%%%%%%%%

\begin{tcolorbox}[colback=blue!5!white, colframe=blue!75!black]
\textbf{Central Reference:} For complete symbol definitions, see
\textbf{Appendix UNMAPPED_GLS} (\textit{Glossary and Notation}).

This section lists only symbols introduced in this appendix.
\end{tcolorbox}

\vspace{0.5em}

\begin{table}[htbp]
\centering
\caption{Symbols Introduced in Appendix UNMAPPED_DTC}
\label{tab:doctype-symbols}
\renewcommand{\arraystretch}{1.3}
\begin{tabular}{@{}clp{6cm}c@{}}
\toprule
\textbf{Symbol} & \textbf{Name} & \textbf{Definition} & \textbf{Also in G?} \\
\midrule
$\mathcal{D}$ & Document Type & 8D vector characterizing a document & NEW \\
$D_1$--$D_8$ & Dimensions & The eight document dimensions & NEW \\
$G_1$--$G_7$ & Goals & The seven communication goals & NEW \\
$\mathcal{S}$ & Structure & Set of emergent structural elements & NEW \\
$\sigma$ & Style & Vector of emergent style parameters & NEW \\
$\mathcal{V}$ & Vocabulary & Collection of word class constraints & NEW \\
$h(\cdot)$ & Structure function & $\mathcal{S} = h(D_1,...,D_8)$ & NEW \\
$s(\cdot)$ & Style function & $\sigma = s(D_1,...,D_8)$ & NEW \\
$v(\cdot)$ & Vocabulary function & $\mathcal{V} = v(\sigma)$ & NEW \\
\bottomrule
\end{tabular}
\end{table}


%%%%%%%%%%%%%%%%%%%%%%%%%%%%%%%%%%%%%%%%%%%%%%%%%%%%%%%%%%%%%%%%%%%%%%%%%%%%%%%

%%%%%%%%%%%%%%%%%%%%%%%%%%%%%%%%%%%%%%%%%%%%%%%%%%%%%%%%%%%%%%%%%%%%%%%%%%%%%%%
\section{Critical Foundations and Objections}
\label{sec:dtc-foundations}
%%%%%%%%%%%%%%%%%%%%%%%%%%%%%%%%%%%%%%%%%%%%%%%%%%%%%%%%%%%%%%%%%%%%%%%%%%%%%%%

\subsection{Objection 1: Scope Limitations}

\textbf{Objection:} The framework presented may have limited applicability
outside the specific domain contexts discussed.

\textbf{Response:} We acknowledge domain-specificity. The framework provides
general principles that should be calibrated to specific contexts. Cross-domain
validation remains an open research question.

\subsection{Objection 2: Measurement Challenges}

\textbf{Objection:} Key constructs may be difficult to measure reliably in
practice.

\textbf{Response:} We recommend using validated scales and instruments where
available, and developing domain-specific proxies where necessary. Sensitivity
analysis should assess robustness to measurement uncertainty.


\section{Open Issues and Research Agenda}
\label{sec:doctype-open}
%%%%%%%%%%%%%%%%%%%%%%%%%%%%%%%%%%%%%%%%%%%%%%%%%%%%%%%%%%%%%%%%%%%%%%%%%%%%%%%

\begin{table}[htbp]
\centering
\caption{Research Roadmap for METHOD-DOCTYPE}
\label{tab:doctype-roadmap}
\renewcommand{\arraystretch}{1.3}
\begin{tabular}{@{}clcl@{}}
\toprule
\textbf{\#} & \textbf{Issue} & \textbf{Priority} & \textbf{Status} \\
\midrule
1 & Validate dimensions empirically & HIGH & Pending \\
2 & Calibrate LLM measurement & HIGH & Pending \\
3 & Extend to non-Western document traditions & MEDIUM & Not started \\
4 & Add temporal dimension (document evolution) & LOW & Conceptual \\
\bottomrule
\end{tabular}
\end{table}


%%%%%%%%%%%%%%%%%%%%%%%%%%%%%%%%%%%%%%%%%%%%%%%%%%%%%%%%%%%%%%%%%%%%%%%%%%%%%%%
\section{References}
\label{sec:doctype-references}
%%%%%%%%%%%%%%%%%%%%%%%%%%%%%%%%%%%%%%%%%%%%%%%%%%%%%%%%%%%%%%%%%%%%%%%%%%%%%%%

\begin{tcolorbox}[colback=blue!5!white, colframe=blue!75!black]
\textbf{Master Bibliography:} For complete reference details, see
\texttt{bcm2\_master\_references.bib} in the repository.
\end{tcolorbox}

\subsection{Key References for This Appendix}

\begin{itemize}[nosep]
    \item \textbf{Register Theory:} Halliday \& Hasan (1989)
    \item \textbf{Genre Analysis:} Swales (1990)
    \item \textbf{Metadiscourse:} Hyland (2005)
    \item \textbf{Readability:} Flesch (1948), Kincaid et al. (1975)
    \item \textbf{Terminology Standards:} ISO 704, ISO 1087
\end{itemize}

\subsection{Appendix-Specific References}

\begin{thebibliography}{99}

\bibitem{halliday1989} Halliday, M.A.K., \& Hasan, R. (1989).
\textit{Language, Context, and Text: Aspects of Language in a Social-Semiotic Perspective}.
Oxford University Press.

\bibitem{swales1990} Swales, J.M. (1990).
\textit{Genre Analysis: English in Academic and Research Settings}.
Cambridge University Press.

\bibitem{hyland2005} Hyland, K. (2005).
\textit{Metadiscourse: Exploring Interaction in Writing}.
Continuum.

\bibitem{flesch1948} Flesch, R. (1948).
A new readability yardstick.
\textit{Journal of Applied Psychology}, 32(3), 221--233.

\end{thebibliography}


%%%%%%%%%%%%%%%%%%%%%%%%%%%%%%%%%%%%%%%%%%%%%%%%%%%%%%%%%%%%%%%%%%%%%%%%%%%%%%%
% CROSS-REFERENCE FOOTER (Required)
%%%%%%%%%%%%%%%%%%%%%%%%%%%%%%%%%%%%%%%%%%%%%%%%%%%%%%%%%%%%%%%%%%%%%%%%%%%%%%%

\vspace{1em}
\hrule
\vspace{0.5em}

\noindent\textit{Cross-references:}
Appendix UNMAPPED_GLS (Central Glossary),
Appendix LLM1 (METHOD-LLMMC),
Appendix CTW (CORE-WHEN),
Appendix UNMAPPED_DTE (REF-DOCTYPE --- theoretical foundations).

\noindent\textit{Master Bibliography:} \texttt{bcm2\_master\_references.bib}


% =============================================================================
% END OF APPENDIX CCC
% =============================================================================
